\documentclass[../../book]{subfiles}

\begin{document}

Finally, we have enough prerequisites to consider one of the main construction blocks of modern zero-knowledge protocols: \textit{commitment schemes}.

\subsection{Basic Definitions and Properties}

The main idea of commitment schemes is specified in the following definition.

\begin{definition}
    A cryptographic commitment scheme allows one party to commit to a chosen statement (such as a value, vector, or polynomial) without revealing the statement itself. The commitment can be revealed in full or in part at a later time, ensuring the integrity and secrecy of the original statement until the moment of disclosure.
\end{definition}

Before delving into the details, here is the high-level intuition of cryptographic commitments.

\begin{intuition}
Imagine putting a letter with some message into a box and locking it with your key. 
You then give that box to your friend, who cannot open it without the key.
In this scenario, you have made a commitment to the message inside the box. 
You cannot change the content of the letter, as it is in your friend's possession. 
At the same time, your friend cannot access the letter since they do not have the key to unlock the box.

\begin{figure}[H]
    \centering
    \includegraphics[width=0.5\linewidth, clip]{contents/images/lecture_5/CommitmentExample.png}

    \caption{Commitment scheme}
\end{figure}
\end{intuition}

From a more formal side, commitments will serve as a short summary of a large pile of data.
For example, suppose you have $n$ scalars $z_1,\dots,z_n \in \mathbb{F}$ (think of $n$ being several millions).
Instead of working with them individually, we can interpolate them into a single polynomial $f \in \mathbb{F}[x]^{\leq n}$ (for example, let $f(\omega_i)=z_i$ for each $i$ for some chosen distinct points $\omega_1,\dots,\omega_n$).
This alone does not give any storage benefits: we still need $n$ scalars to store $f$.
However, using one of the commitment schemes we will consider further, we can build the \emph{commitment} to $f$ that requires only \emph{one} group element to store, regardless of the size of $f$.

Depending on the commitment scheme, we might work with commitments further.
For example, instead of conducting operations over the original data, we can conduct similar operations over committed (succinctly stored) values that will preserve operations over underlying data (the property called \emph{homomorphism}).

Without further ado, let us provide formal definitions.

\begin{definition}[Commitment Scheme]
    Commitment Scheme $\Pi$ is a tuple of two algorithms: $(\mathsf{Setup}, \mathsf{Commit})$.

    \begin{enumerate}
        \item $\mathsf{Setup}(1^{\lambda}) \to \mathsf{pp}$: returns public parameter $\mathsf{pp}$ for both committer (or prover) and verifier;
        \item $\mathsf{Com}(\mathsf{pp}, m) \to c$: returns a commitment $c$ to the message $m$ using public parameters $\mathsf{pp}$;
    \end{enumerate}
\end{definition}

We note that in the literature, one might encounter the third algorithm called $\mathsf{Open}(\mathsf{pp}, m, c) \to \{0,1\}$ that verifies whether given commitment $c$ corresponds to the message $m$.
However, for our applications (with the only exception of \Cref{def:pedersen-poly-commitment}), $\mathsf{Open}(\cdot)$ is trivial:
\begin{equation*}
    \mathsf{Open}(\mathsf{pp}, m, c) = \llbracket \mathsf{Com}(\mathsf{pp}, m) = c \rrbracket,
\end{equation*}
so there is no need to define it separately from $\mathsf{Setup}(\cdot)$ and $\mathsf{Com}(\cdot)$.

\begin{remark}
    In many applications, the commitment scheme also includes the random component $r$ together with a message, which we write as $\mathsf{Com}(m;r)$.
    Such a component is necessary to satisfy the \emph{hiding} property (see \Cref{def:hiding}) of a commitment scheme.
\end{remark}

We now specify security definitions that commitment schemes must satisfy.
The first property is \textbf{hiding} that informally states that a verifier cannot learn any information about the message given only the commitment.
The formal definition is the following.

\begin{definition}[Hiding Property]\label{def:hiding}
    Let $\Pi$ be the commitment scheme. Define the following game for adversary $\mathcal{A}$:
    \begin{algoen}
        \item Challenger and adversary set up public parameters: $\mathsf{pp} \gets \mathsf{Setup}(1^{\lambda})$.
        \item Adversary chooses two messages $m_0,m_1$ and sends to the challenger.
        \item Challenger chooses a random bit $b$, commits to $b$th message by computing\footnote{If randomness is included in the commitment scheme, the challenger additionally passes random $r$ to $\mathsf{Com}(\cdot)$ without revealing it to $\mathcal{A}$.} $c_b \gets \mathsf{Com}(pp, m_b)$, and sends $c_b$ to the adversary.
        \item Adversary guesses a bit $\hat{b}$.
    \end{algoen}

    We define the hiding advantage of $\mathcal{A}$ as $\text{Adv}^{\text{Hide}}_{\mathcal{A},\Pi}(\lambda) := \left| \Pr[b=\hat{b}] - \frac{1}{2} \right|$.
    We say that $\Pi$ is \textbf{hiding} if for any PPT $\mathcal{A}$, $\text{Adv}^{\text{Hide}}_{\mathcal{A},\Pi}(\lambda) = \text{negl}(\lambda)$.
\end{definition}

The second property we require is that the prover cannot find the message $m'$ that opens the commitment $c=\mathsf{Com}(\mathsf{pp}, m)$ without knowing the underlying message $m$.
To put it formally, we give the following definition.

\begin{definition}[Binding Property]
    Let $\Pi$ be the commitment scheme. Define the following game for adversary $\mathcal{A}$:
    \begin{algoen}
        \item Challenger and adversary set up public parameters: $\mathsf{pp} \gets \mathsf{Setup}(1^{\lambda})$.
        \item Adversary chooses $(m_0,m_1,r_0,r_1) \gets \mathcal{A}(\mathsf{pp})$.
    \end{algoen}
        
    Define the advantage in the binding game as:
    \begin{equation*}
        \mathsf{Adv}^{\text{Bind}}_{\mathcal{A}, \Pi}(\lambda) = \Pr[\text{Com}(m_0;r_0) = \text{Com}(m_1;r_1) \; \text{and} \; m_0 \neq m_1].
    \end{equation*}

    We say that the commitment scheme is \textbf{binding} if $\mathsf{Adv}^{\text{Bind}}_{\mathcal{A}, \Pi}(\lambda)=\text{negl}(\lambda)$ for all efficient adversaries $\mathcal{A}$.
\end{definition}

\begin{remark}
    We say that the scheme is \textbf{perfect} binding/hiding if underlying advantages are \emph{exactly zero} for any (not necessarily efficient) adversaries.
\end{remark}

\subsection{Simple Hash-based Commitment Scheme}

We introduce the most basic class of commitment schemes: \emph{hash-based commitments}.
As the name implies, we are going to use a cryptographic hash function \(\mathsf{H}\).
Define the basic hash-based commitment as follows.

\begin{definition}\label{def:hash-based-commitment}
    The \textbf{hash-based commitment} $\Pi_{\mathsf{H}}$ is defined as follows:
    \begin{enumerate}
        \item $\mathsf{Setup}(1^{\lambda})$ initializes the hash function $\mathsf{H}$.
        \item $\mathsf{Com}(\mathsf{pp}, m, r) \to c$ samples the blinding factor $r$ from some finite subset of $\{0,1\}^*$ and outputs $c \gets \mathsf{H}(m \,\|\, r)$.
    \end{enumerate}
\end{definition}

We then claim the following.

\begin{lemma}
    Given a collision-resistant hash function $\mathsf{H}$, the provided $\Pi_{\mathsf{H}}$ above is computationally binding.
\end{lemma}

\begin{proof}
    We proceed using the standard security reduction technique (see \Cref{subsec:security-reduction} if you need a reminder).
    Suppose there exists an efficient adversary $\mathcal{A}$ that breaks the \emph{binding} property of $\Pi_{\mathsf{H}}$.
    However, this implies that the adversary has successfully found a collision ($\mathsf{H}(m_0\,\|\, r_0) = \mathsf{H}(m_1\,\|\, r_1)$ with $m_0 \neq m_1$) that contradicts the fact that $\mathsf{H}$ is collision-resistant.
\end{proof}

Moreover, given a slightly broader background in cryptanalysis, one can show the following property.
\begin{lemma}
    $\Pi_{\mathsf{H}}$ is hiding in the random oracle model for $\mathsf{H}$.
\end{lemma}

\subsection{Pedersen Commitment Scheme}

Pedersen commitments \cite{pedersen1991} allow us to represent arbitrarily large vectors with a single elliptic curve point, while optionally hiding any information about the vector. Pedersen commitment uses a public group $\mathbb{G}$ of order $q$ (written multiplicatively further on) and two random public generators $g$ and $h$ with $h=g^{\rho}$. The secret parameter $\rho$ should be unknown to anyone, otherwise the $\textit{Binding}$ property of the commitment scheme will be violated.
Elliptic curve point $h$ is chosen randomly using ``Nothing-up-my-sleeve'' technique to assure no one knows the discrete logarithm of $h$.

\begin{remark}
    \textbf{Transparent random points generation}

    User can pick the publicly chosen random number (like a hash of a project name or first numbers of $\pi$), and hash that result to obtain a coordinate $x \in \mathbb{Z}_q$. Using a standard try-and-increment strategy, if $x$ does not yield a valid point on the elliptic curve $E(\mathbb{F}_q)$, $x$ is repeatedly incremented until a valid coordinate pair $(x, y)$ is found. Subsequent generators are derived iteratively by hashing the preceding point $(x, y)$ or appending a counter.

    Because the point generation relies on a preimage-resistant cryptographic hash function over a fixed public seed, no party knows the relative discrete logarithms between the generated points.
\end{remark}

Now, we provide the definition of the Pedersen commitment scheme.

\begin{definition}
    Define the \textbf{Pedersen commitment scheme} $\Pi_{\text{Ped}}$ as follows:
    \begin{itemize}
        \item $\mathsf{Setup}(1^{\lambda}) \to \mathsf{pp}$: prover and verifier agree on points $g$ and $h$ on the selected elliptic curve $\mathbb{G}$.
        Let $q:=\#\mathbb{G}=\Theta(2^{\lambda})$.
        \item $\mathsf{Com}(m \in \mathbb{Z}_q; r) \to (c \in \mathbb{G})$: to generate a commitment to $m \in \mathbb{Z}_q$, the prover selects a blinder factor $r \leftarrowS \mathbb{Z}_q$ and outputs $c \gets g^m h^r$.
    \end{itemize}
\end{definition}

We state the following result about the security of the Pedersen commitment scheme.

\begin{lemma}\label{lemma:pedersen-hiding-binding}
    The Pedersen commitment scheme $\Pi_{\text{Ped}}$ is \emph{perfectly} hiding and \emph{computationally} binding under the discrete logarithm assumption (\Cref{def:discrete-log}).
\end{lemma}

\begin{proof}
    \emph{Perfect Hiding.} Notice that $\mathbb{Z}_q \to \mathbb{G}, r \mapsto h^r$ is a bijection. Thus, $h^r$ is uniformly distributed in $\mathbb{G}$.
    Hence, for any fixed $m \in \mathbb{Z}_q$, the commitment $g^mh^r$ is also uniformly distributed in $\mathbb{G}$.

    \emph{Computational Binding.} Assume that there exists an efficient adversary $\mathcal{A}$ that produces two valid openings $(m_0,r_0)$ and $(m_1,r_1)$ with $(m_0,r_0) \neq (m_1,r_1)$ for the same commitment. 
    Hence, $g^{m_0}h^{r_0} = g^{m_1}h^{r_1}$.
    Rearranging yields the identity $h^{r_1-r_0} = g^{m_0-m_1}$. 
    Since $r_0 \neq r_1$, the adversary finds the discrete logarithm of $h$ by providing $(m_0-m_1)(r_1-r_0)^{-1} \in \mathbb{Z}_q$.
\end{proof}

\begin{remark}
    In case the discrete logarithm of $h$ is leaked, the \emph{binding} property can be violated by the \emph{prover}.
    Indeed, notice that the commitment to $(m,r)$ in such case has the form $c = g^{m + r\rho}$.
    Consequently, $(m+\rho, r-1)$ will be another valid yet distinct opening to $c$.
\end{remark}

\subsubsection{Commitments Aggregation}

The Pedersen commitment scheme has a significant advantage compared to hash-based commitments.
Specifically, the Pedersen commitment scheme is \textit{homomorphic}.
Such a property allow us to accumulate multiple commitments into one.
Consider two pairs as an example: $(m_0, r_0)$ and $(m_1, r_1)$.
Suppose we commit to them separately:
\begin{equation*}
    c_0 \gets g^{m_0}h^{r_0}, \quad c_1 \gets g^{m_1}h^{r_1}.
\end{equation*}

Notice that $c_0c_1$ is a commitment to $(m_0+m_1,r_0+r_1)$. Indeed,
\begin{align*}
    c_0c_1 &= g^{m_0}h^{r_0}g^{m_1}h^{r_1} = g^{m_0+m_1}h^{r_0+r_1} = \text{Com}(m_0+m_1;r_0+r_1).
\end{align*}

This works for any (reasonable) number of commitments, so we can encode as many points as we like in a single one.
For example, this is used in \emph{confidential transactions}\footnote{Note, however, that in addition to Pedersen commitments, one needs to attach \textit{range proofs} that underlying numbers are small enough. This is exactly one of the applications of zero-knowledge proofs that we will consider further on.}: given input secret amounts $x_0,\dots,x_n$ and output secret amounts $z_0,\dots,z_m$, we can check condition $x_0+\dots+x_n=z_0+\dots+z_m$ without revealing individual $x_i$'s and $z_i$'s.

This allows to introduce the following definition.

\begin{definition}[Homomorphic Property]
    Let $\Pi=(\mathsf{Setup},\mathsf{Com})$ be a commitment scheme over message space $(\mathcal{M},\times)$ and commitment space $(\mathcal{C},\star)$, both being abelian groups.
    $\Pi$ is called \textbf{homomorphic} if
    \begin{equation*}
        \mathsf{Com}(m_0 \times m_1) = \mathsf{Com}(m_0) \star \mathsf{Com}(m_1) \; \commentgray{for any messages $m_0,m_1 \in \mathcal{M}$.}
    \end{equation*}
\end{definition}

According to this definition, Pedersen commitment scheme is homomorphic (for $\mathcal{M}:=\mathbb{Z}_q$ and $\mathcal{C}:=\mathbb{G}$).

\subsection{Vector Commitment Schemes}

Vector commitment schemes allow to commit to a vector of values rather than a single value and a blinding term.
This is a highly useful tool for us since in zero-knowledge protocols we need to operate over thousands or even millions of variables at once.

\subsubsection{Pedersen Vector Commitment Scheme}

Suppose we have a set of random elliptic curve points $(g_0, \ldots, g_{n-1}, h)$ of a \emph{cyclic} group $\mathbb{G}$ (that we do not know the discrete logarithm of).
Instead of committing to messages in $\mathbb{Z}_q$, we can easily commit to messages in $\mathbb{Z}_q^n$ by slightly modifying Pedersen commitment scheme as follows.

\begin{definition}
    The \textbf{Pedersen vector commitment scheme} $\Pi_{n\text{-Ped}}$ is defined as follows:
    \begin{itemize}
        \item $\mathsf{Setup}(1^{\lambda}) \to \mathsf{pp}$: prover and verifier agree on the cyclic elliptic curve $\mathbb{G}$ of order $q:=\#\mathbb{G}=\Theta(2^{\lambda})$.
        They select $g_0,\dots,g_{n-1},h \leftarrowS \mathbb{G}$.
        \item $\mathsf{Com}(\mathbf{m} \in \mathbb{Z}_q^n; r \in \mathbb{Z}_q) \to (c \in \mathbb{G})$: to commit to $\mathbf{m} \in \mathbb{Z}_q^n$, the prover selects a blinder factor $r \leftarrowS \mathbb{Z}_q$ and outputs
        \begin{equation*}
            c \gets h^r\prod_{i \in [n]} g_i^{m_i}.
        \end{equation*}
    \end{itemize}
\end{definition}

Similarly to \Cref{lemma:pedersen-hiding-binding}, one can demonstrate the following.

\begin{lemma}
    The Pedersen vector commitment scheme $\Pi_{n\text{-Ped}}$ is \emph{perfectly} hiding and \emph{computationally} binding under the discrete logarithm assumption.
\end{lemma}

\subsubsection{Merkle Tree Commitment Scheme}

A naïve approach for a vector commitment would be to hash the whole vector. 
A more sophisticated scheme uses divide-and-conquer approach by building a binary tree out of vector elements: consider \Cref{fig:merkle-tree-structure}.

\begin{figure}[H]
    \centering\includegraphics[width=0.9\linewidth, clip]{contents/images/lecture_5/MerkleTree.png}

    \caption{Merkle Tree structure}\label{fig:merkle-tree-structure}
\end{figure}

A \emph{Merkle Tree} is a data structure to efficiently and securely verify the commitments to a vector of data. 
It is a binary tree where each leaf node represents a hash of a data block, and each non-leaf node is a hash of its child nodes' concatenated hashes. 
The top node, called the \emph{root hash} or \emph{Merkle root}, uniquely represents the entire data set. 
By comparing this root with a known valid root, one can quickly verify the authenticity and integrity of the data without needing to examine the entire dataset. 

To prove the inclusion of an element in the tree, a corresponding Merkle Branch is used. In the example below (\Cref{fig:merkle-tree-path}), $M_1$ inclusion is proved, and 
$(M_2, \mathsf{H}(M_3 \parallel M_4), \mathsf{H}(\mathsf{H}(M_5 \parallel M_6) \parallel \mathsf{H}(M_7 \parallel M_8)))$ is an inclusion branch vector.

\begin{figure}[H]
    \centering\includegraphics[width=0.9\linewidth, clip]{contents/images/lecture_5/MerkleTreeProof.png}

    \caption{Merkle Tree inclusion proof branch}\label{fig:merkle-tree-path}
\end{figure}

One of the key advantages of a Merkle tree is that it allows for the selective disclosure of specific elements within the data set without revealing the rest. 

\subsubsection{Pedersen Polynomial Commitment Scheme}

\textbf{Polynomial commitment schemes} can be used to prove that the committed polynomial satisfies certain properties (most commonly, an evaluation condition $f(\rho) = y$ for a randomly chosen point $\rho$ and that $\deg f < d$), without revealing what the polynomial is.
The commitment is generally succinct, which means that it is much smaller than the polynomial it represents. Typically, polynomial commitment scheme is a heart of any complex \textit{zero-knowledge} proof system.

\begin{remark}
    Before delving into polynomial commitment schemes, we make an essential remark.
    The subsequent scheme will be formulated in the form ``prover sends $\mathsf{A}$ to the verifier, verifier responds with $\mathsf{B}$, prover responds with $\mathsf{C}$''.
    Such formulations are \emph{very} common in the zero-knowledge field.

    This might be confusing given that previously considered protocols are \emph{non-interactive}.
    However, using tools from \Cref{section:intro-zk} (specifically, \Cref{subsection:fiat-shamir}), we will be able to make them non-interactive rather straightforwardly.
\end{remark}

\textbf{Pedersen Polynomial Commitment Scheme.} Suppose the prover wants to convince the verifier that he knows evaluation of a polynomial $f \in \mathbb{F}_{q}[T]^{<d}$ at some point $\rho \in \mathbb{F}_q$ (\textit{i.e.}, $f(\rho) = y$) without revealing anything about $f$. Let $\mathbb{G}$ be a cyclic group of order $q$ where the \textit{discrete logarithm} problem is hard. 
Let $g,h \in \mathbb{G}$ be points with unknown discrete logarithms between them. Then, we define the basic polynomial commitment scheme as follows.

\begin{definition}\label{def:pedersen-poly-commitment}
    The \textbf{Pedersen polynomial commitment scheme} $\Pi_{\text{poly}}$ is defined as follows:
    \begin{enumerate}
        \item $\mathsf{Setup}(1^{\lambda}) \to \mathsf{pp}$: the prover and verifier agree on a cyclic group $\mathbb{G}$ of order $q = \Theta(2^{\lambda})$ and two random points $g,h \leftarrowS \mathbb{G}$.
        \item $\mathsf{Com}(f \in \mathbb{F}_q[T]^{<d}) \to ((y,z) \in \mathbb{Z}_q^2, \mathbf{c} \in \mathbb{G}^d)$: given $f=\sum_{i \in [d]}f_iT^i$, the prover and verifier interact as follows:
        \begin{algoen}
            \item the prover selects random $\bm{\gamma} \leftarrowS \mathbb{F}_q^d$ and computes $d$ Pedersen commitments $c_i \gets g^{f_i}h^{\gamma_i}$ for $i \in [d]$;
            \item the verifier samples $\rho \leftarrowS \mathbb{F}_q$ and sends it to prover;
            \item the prover constructs the blinded polynomial $\gamma = \sum_{i \in [d]}\gamma_i T^i$ and outputs $(y,z) \gets (f(\rho), \gamma(\rho))$.
        \end{algoen}
        \item $\mathsf{Open}(\mathbf{c}, \rho, y, z) \to \{0,1\}$: the verifier outputs $1$ if and only if the following check passes: $\prod_{i \in [d]} c_i^{\rho^i} = g^yh^z$.
    \end{enumerate}
\end{definition}

We provide the following statement to prove security of the aforementioned scheme.

\begin{lemma}
    The Pedersen polynomial commitment scheme from \Cref{def:pedersen-poly-commitment} is perfectly hiding and computationally binding based on discrete logarithm assumption.
\end{lemma}

\begin{proof}
    Despite the seemingly high complexity of the protocol, the proof is similar to \Cref{lemma:pedersen-hiding-binding}.

    \textit{Perfect hiding.} To show that the scheme is perfectly hiding, we need to prove that the output distribution $(\mathbf{c},y,z)$ is independent of $f$.
    Specifically, one can show that we can \emph{simulate} it as follows.

    \underline{$\mathsf{Sim}(\rho,y) \to (\mathbf{c},z):$}
    \begin{algoen}
        \item sample $z \leftarrowS \mathbb{F}_q$;
        \item pick any $j \in [d]$ such that $\rho^j \neq 0$;
        \item sample $c_0,\dots,c_{j-1},c_{j+1},\dots,c_{d-1} \leftarrowS \mathbb{G}$;
        \item set $c_j \gets \left(g^yh^z/\prod_{i \neq j}c_i^{\rho^i}\right)^{\rho^{-j}}$;
        \item \textbf{output} $(\mathbf{c},z)$.
    \end{algoen}

    \begin{exercise}
        Show that the distribution of $(\mathbf{c},z)$ from the Pedersen commitment scheme $\Pi_{\text{poly}}$ (\Cref{def:pedersen-poly-commitment}) is exactly $\mathsf{Sim}(\rho,y)$.
    \end{exercise}

    \textit{Computational binding.} The proof is exactly as in \Cref{lemma:pedersen-hiding-binding}.
    Notice that if some efficient adversary finds two valid openings $(y,z) \neq (y',z')$, then 
    \begin{equation*}
        \prod_{i \in [d]} c_i^{\rho^i} = g^yh^z = g^{y'}h^{z'}.
    \end{equation*}
    Consequently, the adversary solves discrete logarithm for $h$ by submitting $(y-y')(z'-z)^{-1}$ to the discrete logarithm game.
    This is a contradiction.
\end{proof}

\begin{remark}
    The idea of a simulator used in the proof might be confusing at this point.
    For more details, refer to \Cref{subsection:simulation}.
\end{remark}

\subsubsection{KZG Commitment Scheme}

Unfortunately, the scheme from the previous section is not succinct, as the size of the commitment grows linearly with the degree of the polynomial.
However, using pairing (\Cref{section:ecpairing}), one can build the commitment scheme that requires only a single point in $\mathbb{G}$ to commit to the whole polynomial.
This is exactly the KZG scheme that we describe in this section.

\textbf{Basic idea.} As before, suppose the prover wants to convince the verifier that $f(\rho)=y$ for some random point $\rho \leftarrowS \mathbb{F}_q$ chosen by the verifier.
Now we employ \Cref{th:root-div}. Notice that $\rho$ is a zero of a polynomial $f(T)-y$. 
Consequently, there exists some polynomial $g \in \mathbb{F}_q[T]^{<d}$ (commonly referred to as a \emph{quotient polynomial}) such that $f(T)-y = (T-\rho)g(T)$.

To turn this into a sound scheme, before receiving $\rho$, the prover must first \emph{commit} to polynomial $f$ (call it $\mathsf{com}(f)$).
After obtaining $\rho$, the prover commits to $q$ (call it $\mathsf{com}(q)$) and then the verifier must check the identity $f(T)-y = (T-\rho)g(T)$.
This check is most efficiently done using \textit{Schwartz-Zippel lemma} (\Cref{lemma:one-sz}).
Specifically, the verifier samples another challenge $\tau \leftarrowS \mathbb{F}_q$ and checks whether $f(\tau)-y=(\tau-\rho)g(\tau)$ using $\mathsf{com}(f)$ and $\mathsf{com}(g)$.
However, how exactly can we commit to these polynomials?

\textbf{Committing to polynomials.} Let $\mathbb{G}_1$ and $\mathbb{G}_2$ be two cyclic groups with corresponding generators $g_1,g_2$.
Assume that we have a secure and efficient pairing $e: \mathbb{G}_1 \times \mathbb{G}_2 \to \mathbb{G}_T$.
To commit to a polynomial $f \in \mathbb{F}_q[T]^{<d}$, we compute $\mathsf{com}(f) \gets g_1^{f(\tau)}$.
Then, using two commitments $\mathsf{com}(f)$ and $\mathsf{com}(g)$, we can do the main using
\begin{equation*}
    e(\mathsf{com}(f)/g_1^y, g_2) = e(\mathsf{com}(g), g_2^{\tau-\rho}).
\end{equation*}

However, notice the following crucial problem: the commitment uses point $\tau$ \emph{before} the verifier sends the challenge $\rho$.
More specifically, the knowledge of $\tau$ breaks the binding property of a commitment to a polynomial.
Fortunately, the prover can commit to $f$ \emph{even without knowing $\tau$}.

\textbf{Main construction.} Let us write down $g_1^{f(\tau)}$ coefficient-wise:
\begin{equation*}
    g_1^{f(\tau)} = g_1^{\sum_{i \in [d]}f_i\tau^i} = \prod_{i \in [d]}g_1^{f_i\tau^i} = \prod_{i \in [d]}(g_1^{\tau^i})^{f_i}.
\end{equation*}

Now observe the following: we can compute $g_1^{f(\tau)}$ solely based on the set $\mathsf{pp} := (g_1,g_1^{\tau},g_1^{\tau^2},\dots,g_1^{\tau^{d-1}})$.
Moreover, given that discrete logarithm assumption holds in $\mathbb{G}_1$, the set $\mathsf{pp}$ does not allow to restore $\tau$.
Consequently, the trusted party can generate the set $\mathsf{pp}$ once, \textbf{delete $\tau$}, and publish $\mathsf{pp}$.
Then anyone can compute commitments to polynomials using $\mathsf{pp}$ and do the check described above.

\begin{remark}
    We stress once more that KZG commitment scheme requires the so-called \textbf{trusted setup} which consists of generating the set $\{g_1^{\tau^i}\}_{i \in [d]}$ by a trusted party and then deleting $\tau$.
    This is arguably the largest concern with employing KZG in practice.
    Fortunately, using appropriate distributed generation process (informally called a \emph{$\tau$-ceremony}), one can generate this set safely given \textit{at least one} honest participant.
\end{remark}

This way, we give the following definition of the KZG commitment.

\begin{definition}\label{def:kzg-com}
    The \textbf{KZG commitment scheme} $\Pi_{\text{KZG}}$ is defined as follows:
    \begin{itemize}
        \item $\mathsf{Setup}(1^{\lambda}) \to \mathsf{pp}$ consists of two steps: 
        \begin{algoen}
            \item the prover and verifier agree on the group $\mathbb{G}$ of order $q := \Theta(2^{\lambda})$;
            \item the trusted party chooses a \emph{toxic parameter} $\tau \leftarrowS \mathbb{Z}_q$ and outputs $\mathsf{pp} \gets \{g_1^{\tau^i}\}_{i \in [d]}$. It then \emph{deletes $\tau$};
        \end{algoen}
        \item $\mathsf{Com}(f \in \mathbb{F}_q^{<d}[T]) \to \mathbb{G}$: given a polynomial $f = \sum_i f_iT^i$, the prover parses $\{h_i\}_{i \in [d]} \gets \mathsf{pp}$ and outputs $\prod_{i \in [d]} h_i^{f_i}$.
    \end{itemize}
\end{definition}

\begin{lemma}
    The KZG commitment scheme is computationally binding given $d$-power Diffie-Hellman assumption\footnote{While we have not defined this assumption rigorously, informally it asks to compute $g^{\tau^{d+1}}$ given the set $\{g^{\tau^i}\}_{i\leq d}$ (for $\tau \leftarrowS \mathbb{F}_q$).}.
\end{lemma}

\begin{remark}
    As stated, the KZG scheme is \emph{not hiding}. To make it hiding, one needs to introduce a blinding factor $r \leftarrowS \mathbb{Z}_q$ and form the commitment as
    \begin{equation*}
        h^r\prod_{i \in [d]}(g^{\tau^i})^{f_i} \quad \commentgray{where $h \in \mathbb{G}$.}
    \end{equation*}
\end{remark}

Consequently, we build the following evaluation proof protocol.

\begin{definition}[KZG Evaluation Proof]\label{def:kzg-proof}
    Suppose prover wants to show that $f(\rho)=y$ for a random point $\rho \leftarrowS \mathbb{F}_q$.
    The protocol proceeds as follows.
    \begin{itemize}
        \item $\mathsf{Setup}(1^{\lambda}) \to \mathsf{pp}$: this method consists of two steps:
        \begin{algoen}
            \item prover and verifier agree on groups $\mathbb{G}_1$ and $\mathbb{G}_2$ with efficiently computable pairing $e: \mathbb{G}_1 \times \mathbb{G}_2 \to \mathbb{G}_T$.
            Suppose the order of $\mathbb{G}_1$ is $q=\Theta(2^{\lambda})$;
            \item prover and verifier fix KZG commitment scheme $\Pi_{\text{KZG}}$ over $\mathbb{G}_1$ and append its public parameters to $\mathsf{pp}$;
            \item while generating trusted setup, the trusted party adds $g_2^{\tau}$ to $\mathsf{pp}$.
        \end{algoen}
        \item $\mathsf{Prove}(\mathsf{pp},f,\rho,y) \to \pi$: the prover conducts the following steps: 
        \begin{algoen}
            \item compute $g(T) \gets (f(T)-y)/(T-\rho)$ over $\mathbb{F}_q[T]$;
            \item \textbf{output} $\pi \gets (\Pi_{\text{KZG}}.\mathsf{Com}(f), \Pi_{\text{KZG}}.\mathsf{Com}(g))$.
        \end{algoen}
        \item $\mathsf{Verify}(\pi, \rho, y)$: verifier parses $(\mathsf{com}(f), \mathsf{com}(g)) \gets \pi$ and checks
        \begin{equation*}
            e(\mathsf{com}(f)/g_1^y, g_2) = e(\mathsf{com}(g), g_2^{\tau}/g_2^{\rho}).
        \end{equation*}
    \end{itemize}
\end{definition}

\subsection{Exercises}

\vspace{-1mm}

\begin{xexercise}
    {Exercise 1.}
    {Dmytro and Denis were watching a horse race. Confident in
    his ability to predict the outcome, Dmytro decided to commit to his
    prediction. However, in his haste, he forgot to use a blinding factor. Now,
    Dmytro is concerned that Denis might discover his prediction before the race
    ends, which would defeat the purpose of his commitment.
    
    We define a dummy hash function $\mathsf{H}: \mathbb{Z} \to \mathbb{F}_{41}$ as $\mathsf{H}(x)
    = 13x+17$. Dmytro used a \textit{hash-based commitment} (\Cref{def:hash-based-commitment}) and $\mathsf{H}$ as a hash
    function. The set of race horse numbers is $\{3, 5, 8, 15\}$. Help Denis to
    find out the horse number Dmytro has made a commitment to, if the
    commitment equals $39$.}
    {4}
    {
        \item $3$.
        \item $5$.
        \item $8$.
        \item $15$.
    }
\end{xexercise}

\begin{xexercise}
    {Exercise 2.}
    {Denis made a setup (points $g$ and $h$) for a Pedersen
    commitment scheme and commited values $(m_0, r_0) = (3, 7)$ to Dmytro by sending him
    $c = g^3h^7$. Dmytro did not verify the setup. Turns out that Denis knows
    that $h=g^6$. Denis is planning to send a different message from the one he
    originally committed to: $m_1 = 15$. Which values of $(m_1, r_1)$ should he send
    to Dmytro for them to be a valid opening to $c$ as well?}
    {4}
    {
        \item $(15, 5)$
        \item $(15, 7)$
        \item $(15, 4)$
        \item $(3, 5)$
    }
\end{xexercise}

\textbf{Exercise 3.} We define a dummy hash function $\mathsf{H}: \mathbb{Z} \times
\mathbb{Z} \to \mathbb{F}_{41}$ as $\mathsf{H}(x, y) = 3x+7y$. You have a Merkle tree
built with depth $4$ using the hash function $\mathsf{H}$ with a root being $37$.
Which inclusion proof is valid for element $3$? Position defines how leaves
should be hashed:
\begin{itemize}
    \item if $\ell$, then $h_i \gets \mathsf{H}(h_{i-1}, \mathsf{branch}[i])$,
    \item if $r$, then $h_i \gets \mathsf{H}(\mathsf{branch}[i], h_{i-1})$.
\end{itemize}
\begin{enumerate}[(A)]
    \item Branch: $[4, 16, 13]$, Position: $[\ell,r,\ell]$
    \item Branch: $[1, 40, 3]$, Position: $[\ell, \ell, \ell]$
    \item Branch: $[5, 12, 13]$, Position: $[r, r, \ell]$
    \item Branch: $[4, 17, 13]$, Position: $[\ell, r, \ell]$
\end{enumerate}

\textbf{Exercise 4.} Given a polynomial $f(T) = T^3 - 10T^2 + 31T - 30$, Oleksandr wants to prove that $p(2)=0$. 
To do that, according to the KZG commitment scheme (\Cref{def:kzg-proof}), he constructs the quotient polynomial $g(T)$ and wants 
to show that $f(\tau) \cdot (\tau - 2) = g(\tau)$. Assuming Oleksandr has conducted these steps correctly, 
what value of $g(T)$ has Oleksandr calculated?

\vspace{-3mm}

\begin{multicols}{2}
    \begin{enumerate}[(A)]
        \item $g(T) = 2T^2 + 4T - 6$
        \item $g(T) = T^3 - 10T^2 + 30T - 28$
        \item $g(T) = T^2 - 8T + 15$
        \item $g(T) = T^2 + 5T + 18$
    \end{enumerate}
\end{multicols}

\end{document}